\section{Direct stopping for the original economic decision}
\label{supp:decisionstopping126}

This section establishes the stopping and replacement rules used by the
new direct service.  It assembles the inherited original-support results
into a complete service guarantee; it does not claim new convexity or
root-finding principles.  All statements condition on valid observations
of the same original Bellman objective.  The critic supplies a feasible
proposal and never supplies an unverified derivative of that objective.

Fix one state and the original risk law.  Let $A>0$, and let
$q:[0,A]\to\mathbb R$ be continuously differentiable, $\mu$-strongly
convex and have $L$-Lipschitz derivative, with $0<\mu\le L$ and
$|q'|\le M$.  Its unique constrained minimizer is $a^*$ and its value is
$q^*=q(a^*)$.  The complete policy with the unchanged future controller
has value
\begin{equation}\label{eq:directfuture126}
 J(a)=q(a)+\eta(a),\qquad 0\le\eta(a)\le\delta.
\end{equation}
No continuity or differentiability of $\eta$ is assumed.  The comparison
objective is $q^*$, the original optimum with optimal future decisions;
we do not replace it by $\min J$.

Let the retained original observation collection be finite and nonempty,
as it is after the direct service has bought its first support.  Every
observation $j$ has an action $a_j$, an interval
$W_j\le q(a_j)\le U_j$, and a derivative interval
$g_j^-\le q'(a_j)\le g_j^+$.  Its global support supplies a single
$p_j\in[p_j^-,p_j^+]$ such that
\[
 q(t)\ge W_j+p_j(t-a_j)+\frac\mu2(t-a_j)^2
 \quad(0\le t\le A).
\]
The support and derivative intervals are obtained from the original
continuous-risk calculation and the finite-radius conversion of
Theorem~\ref{thm:supportgradient122}.  Neither an approximate actor nor
a fitted continuation derivative establishes these premises.  We also
retain any independently valid original-optimum bound $B_0\le q^*$.

\subsection{Transported derivatives and the constrained optimizer}

For each observation define the continuous piecewise-affine functions
\begin{align*}
 \lambda_j(t)&=g_j^-+
 \begin{cases}L(t-a_j),&t<a_j,\\ \mu(t-a_j),&t\ge a_j,\end{cases}\\
 \upsilon_j(t)&=g_j^++
 \begin{cases}\mu(t-a_j),&t<a_j,\\ L(t-a_j),&t\ge a_j.\end{cases}
\end{align*}
The direct service uses
$\lambda(t)=\max_j\lambda_j(t)$ and
$\upsilon(t)=\min_j\upsilon_j(t)$.  The global cap $M$ remains
available for the separate distance bound below.  For any continuous
strictly increasing $f$ define its constrained zero by
\[
 Z(f)=\begin{cases}
 0,&f(0)\ge0,\\
 A,&f(A)\le0,\\
 \text{the unique zero of }f\text{ in }(0,A),&\text{otherwise}.
 \end{cases}
\]

\begin{lemma}[A nested optimizer enclosure]
\label{lem:directbracket126}
The inequalities $\lambda\le q'\le\upsilon$ hold on $[0,A]$.
Both envelopes are continuous and satisfy
$f(t)-f(s)\ge\mu(t-s)$ for $t>s$.  Consequently
\begin{equation}\label{eq:directbracket126}
 \ell=Z(\upsilon)\le a^*\le Z(\lambda)=r.
\end{equation}
Retaining all old observations when adding a valid new observation can
only increase $\ell$ and decrease $r$.  If
\begin{equation}\label{eq:stationary126}
 \upsilon(0)\le0\le\lambda(A),
\end{equation}
then $q'(a^*)=0$.  This condition allows a boundary optimum whose
one-sided derivative is zero; a general boundary optimum need not be
stationary.
\end{lemma}

\begin{proof}
For $t\ge a_j$, strong convexity and Lipschitz continuity of $q'$
give
$\mu(t-a_j)\le q'(t)-q'(a_j)\le L(t-a_j)$.
For $t<a_j$ the inequalities reverse the two curvature coefficients
when multiplied by the negative displacement.  This proves
$\lambda_j(t)\le q'(t)\le\upsilon_j(t)$, hence the asserted envelope
inequalities.  Each transported function is continuous and its slopes
are $\mu$ or $L$.  If every $f_j$ satisfies
$f_j(t)\ge f_j(s)+\mu(t-s)$, its finite maximum and finite minimum
satisfy the same inequality.  Thus both envelopes are continuous and
strictly increasing.

For such functions, pointwise order reverses the order of constrained
zeros: $f\le g$ implies $Z(g)\le Z(f)$.  To verify this, an interior
zero of $g$ has $f\le0$, so it cannot lie to the right of $Z(f)$;
if the constrained zero is an endpoint, the relevant endpoint sign
gives the same conclusion.  The first-order variational inequality
for minimizing $q$ on $[0,A]$ gives $a^*=Z(q')$, proving
\eqref{eq:directbracket126}.  Additional observations increase
$\lambda$ and decrease $\upsilon$ pointwise, so their constrained
zeros give the stated nesting.  Finally \eqref{eq:stationary126}
implies $q'(0)\le0\le q'(A)$.  By continuity and strict monotonicity,
$q'$ has exactly one zero in $[0,A]$, which is its constrained
minimizer.  Without these signs a constrained minimum can have a
nonzero outward derivative.
\end{proof}

\subsection{Projected gaps and separate stopping gates}

For a feasible realized action $a$ and a derivative interval
$g_-\le q'(a)\le g_+$, put
\begin{align}
 d(z)&=\operatorname{proj}_{[-a,A-a]}(-z/\mu),\nonumber\\
 D_\mu(a;g_-,g_+)&=
 \max_{z\in\{g_-,g_+\}}
       \left\{-z\,d(z)-\frac\mu2d(z)^2\right\}.
 \label{eq:projectedgap126}
\end{align}
The maximum is nonnegative because the feasible displacement set
contains zero.  This is the projected strong-convexity calculation
already used by the inherited action service.  In particular, it
recognizes a constrained optimum without requiring its unconstrained
gradient to vanish.

For completeness, write $C^-(u,v)\le q(v)-q(u)\le C^+(u,v)$ for
the integrals of the appropriate transported derivative envelopes,
with their order reversed when $v<u$ and both equal to zero when
$u=v$.  Define
\begin{align*}
 S_j^-(t)&=W_j+\min\{p_j^-(t-a_j),p_j^+(t-a_j)\}
                         +\frac\mu2(t-a_j)^2,\\
 L_q(a)&=\max_j\{S_j^-(a),\ W_j+C^-(a_j,a)\},\\
 U_q(a)&=\min_j\{U_j+C^+(a_j,a)\},\\
 B&=\max\left\{B_0,\ \min_{t\in[0,A]}\max_j S_j^-(t)\right\}.
\end{align*}
The finite piecewise-quadratic minimum is evaluated by retaining all
breakpoints and feasible stationary points.  It does not assume that
the interval-slope support envelope is convex.  These are the
original-level quantities used by the direct implementation.

\begin{proposition}[A complete direct action certificate]
\label{prop:directcertificate126}
Let $[\ell,r]$ be the bracket in Lemma~\ref{lem:directbracket126},
let $\rho(a)=\max\{|a-\ell|,|a-r|\}$, and use
$g_-=\lambda(a)$ and $g_+=\upsilon(a)$ in
\eqref{eq:projectedgap126}.  Each of the following is an upper bound
on $J(a)-q^*$:
\begin{equation}\label{eq:directroutes126}
 U_q(a)+\delta-B,\qquad
 \delta+D_\mu(a;g_-,g_+),\qquad
 \delta+M\rho(a).
\end{equation}
When \eqref{eq:stationary126} is certified,
$\delta+L\rho(a)^2/2$ is a further upper bound.  Their minimum,
$G(a)$, is therefore a valid nonnegative original-policy regret bound.
Separately,
\begin{equation}\label{eq:directlevel126}
 L_q(a)\le J(a)\le U_q(a)+\delta,
 \qquad W_J(a)=U_q(a)+\delta-L_q(a).
\end{equation}
Hence the two tests $G(a)\le\tau$ and $W_J(a)\le\omega$ establish
the declared original-regret and absolute-value-width requirements.
A smaller prescribed target $\tau_*\le\tau$ is an additional
regret outcome and requires $G(a)\le\tau_*$ separately.

For a fixed action, retaining valid old observations makes every
available bound and the exact width nonincreasing.  A stationary
certificate, once established, remains established.  Thus an already
qualified unchanged action stays qualified under exact accumulation
of valid information.  Moreover
$|a-a^*|\le\sqrt{2G(a)/\mu}$; this quantity is a bound, not a
measurement of realized action error.
\end{proposition}

\begin{proof}
Strong convexity gives
\[
 q(a+d)\ge q(a)+q'(a)d+\frac\mu2d^2
 \quad(-a\le d\le A-a).
\]
Consequently
\[
 q(a)-q^*\le
 \max_{-a\le d\le A-a}
          \{-q'(a)d-\mu d^2/2\}.
\]
For a fixed $z$ the maximizer is exactly $d(z)$.  The maximum over
displacements is a convex function of $z$, being a supremum of
affine functions of $z$.  Its maximum on $[g_-,g_+]$ therefore
occurs at an endpoint, proving the projected bound.  The argument
includes $a=0,A$.  For example, if $a=A$ and $g_+\le0$, each
projected displacement is zero and the projected gap vanishes.

The bracket gives $|a-a^*|\le\rho(a)$.  Integrating the derivative
cap yields $q(a)-q^*\le M\rho(a)$.  If $q'(a^*)=0$, the smooth
upper Taylor inequality gives
$q(a)-q^*\le L|a-a^*|^2/2\le L\rho(a)^2/2$ instead.
Adding $0\le\eta(a)\le\delta$ proves these direct regret routes.

Every raw support $S_j^-$ is below $q$, since its minimum over the
reported slope interval is no larger than the value at the true
shared slope.  Integrating valid derivative envelopes also gives
valid contrasts.  It follows that
$L_q(a)\le q(a)\le U_q(a)$ and $B\le q^*$.
Equation~\eqref{eq:directfuture126} now proves the level gap and
\eqref{eq:directlevel126}.  Taking the minimum of valid upper bounds
preserves validity.  No error is differentiated and the future
allowance is added once.

Adding observations tightens both contrast endpoints, increases each
maximum defining a lower bound, and decreases each minimum defining
an upper bound.  It increases the maximum of the raw supports
pointwise, and hence cannot decrease its minimum.  Thus $B$ and
$L_q(a)$ cannot decrease, and $U_q(a)$ cannot increase.  The derivative
interval narrows, so its robust projected gap cannot increase.
The bracket nests, so $\rho(a)$ cannot increase.
The endpoint signs in \eqref{eq:stationary126} are preserved because
$\upsilon$ decreases and $\lambda$ increases.  These observations
prove monotonicity of the bounds, their minimum and the exact width.
Finally the constrained minimizer satisfies
$q'(a^*)(a-a^*)\ge0$.  Strong convexity then gives
$q(a)-q^*\ge\mu|a-a^*|^2/2$.  Since $q(a)\le J(a)$, the
claimed distance bound follows.
\end{proof}

The implementation converts $L_q$ downwards, $U_q+\delta$ and $G$
upwards to finite binary64 endpoints, and rounds their reported width
upwards.  It checks both the exact rational tests and the reported
outward tests.  A claimed pass therefore requires an outward regret
no larger than the stated tolerance and an outward width no larger
than its separate tolerance.  The mathematical guarantee concerns
these outward enclosures; a nonfinite or inconsistent numerical
conversion is not a successful certificate.

\begin{proposition}[Why the two gates cannot be merged]
\label{prop:separategates126}
An action-regret certificate alone need not determine an absolute
policy-value interval or an original-optimum lower bound.  In
particular, $J^+(a)-G(a)$ need not be a lower bound for $q^*$ when
$J^+(a)$ is merely an upper enclosure and $G(a)$ is a direct regret
bound.  Conversely, exact knowledge of $J(a)$ does not certify a
small decision regret without information about alternatives.
\end{proposition}

\begin{proof}
On $[0,1]$, consider the compatible family
$q_C(t)=C+(t-1/2)^2$, $0\le C\le V$, for a rational $V>0$.
Its derivative and unique optimizer $a=1/2$ are known exactly.  The
common support $(t-1/2)^2$ is valid for every member, while the
absolute-level interval at the optimizer is $[0,V]$.  Under
$0\le\eta\le\delta$, the direct regret bound is $G(a)=\delta$
and the upper policy endpoint is $J^+(a)=V+\delta$.  For $C=0$
and $\eta\equiv0$, the true optimum is zero whereas
$J^+(a)-G(a)=V>0$.  An arbitrarily accurate optimizer certificate
therefore cannot remove the unresolved additive level.  It follows
only that $q^*\ge J(a)-G(a)\ge J^-(a)-G(a)$, using a lower,
not an upper, policy endpoint.  For the converse, take
$q(t)=(t-1)^2$, $\eta=0$, and $a=0$.  Its value at the chosen
action is exactly one, with zero value uncertainty, but its regret
is one.  Both examples satisfy the stated smooth strong-convexity
assumptions.
\end{proof}

\subsection{Strict policy changes and rounded localization}

\begin{proposition}[Same-objective replacement]
\label{prop:strictchange126}
Let $a$ be a completed qualified action and let a distinct realized
candidate $b$ satisfy both qualification gates.  Set
\[
 u(a,b)=\min\{C^+(a,b),\ U_q(b)-L_q(a)\}.
\]
If $u(a,b)+\delta<0$, replacing $a$ by $b$ decreases the original
complete-policy objective by at least
$-u(a,b)-\delta>0$.  No claim of such improvement follows merely
from a smaller regret bound.  For $b=a$, the actual complete-policy
contrast is exactly zero under the unchanged future controller.
\end{proposition}

\begin{proof}
Both entries in the displayed minimum bound $q(b)-q(a)$ from
above.  Hence
\[
 J(b)-J(a)=q(b)-q(a)+\eta(b)-\eta(a)
          \le u(a,b)+\delta.
\]
The result follows.  The single allowance is sufficient because
$\eta(b)\le\delta$ and $\eta(a)\ge0$.  Without further regularity
or coupling of the future errors it cannot in general be removed:
the admissible assignments $\eta(b)=\delta$, $\eta(a)=0$ attain
that allowance.  Identical actions use the identical complete
future policy, so their actual values are equal.  The first
qualified action requires no comparison with an earlier completed
policy; its selection establishes qualification only.
\end{proof}

\begin{lemma}[The actual rounded-midpoint condition]
\label{lem:roundstop126}
Let a valid old optimizer interval be $I=[\ell,r]$, of width $w$,
and let a feasible queried action $a$ have
$\epsilon=|a-(\ell+r)/2|$.  Suppose its derivative interval is
$[g_-,g_+]$.  Intersect $I$ with $[0,a]$ if $g_->0$, with
$[a,A]$ if $g_+<0$, and otherwise with
$[a-g_+/\mu,a-g_-/\mu]$.  The resulting interval remains valid
and has width at most
\begin{equation}\label{eq:roundedprogress126}
 \max\{w/2+\epsilon,\ (g_+-g_-)/\mu\}.
\end{equation}
When $I$ is the bracket of the retained old envelopes, the new full
transported-envelope bracket is at least as informative as this update.

For lattice spacing $h>0$, the implemented feasible downward rounding
$a=h\lfloor ((\ell+r)/2)/h\rfloor$ satisfies
$0\le(\ell+r)/2-a<h$.  Whenever the actual updated bracket passes
the checked inequality $w_{\rm new}\le w/2+\epsilon$, a genuine
localization improvement with that rounding allowance has been
certified.  After $m$ consecutive such checked updates from width
$w_0$,
\[
 w_m\le 2^{-m}w_0+2h(1-2^{-m}).
\]
These statements do not assert that a finite prescribed precision
ladder must produce a passing update.
\end{lemma}

\begin{proof}
If $q'(a)>0$, monotonicity and the constrained optimality condition
imply $a^*\le a$; if $q'(a)<0$, they imply $a^*\ge a$.  This
proves the two sign-resolved intersections, whose widths are at most
$w/2+\epsilon$.  If $g_-\le0\le g_+$ and $a^*<a$, strong
monotonicity and $q'(a^*)\ge0$ when the minimizer is the left
boundary give $q'(a)\ge\mu(a-a^*)$.  Thus
$a^*\ge a-g_+/\mu$, while the proposed upper endpoint is at
least $a$.  If $a^*>a$, the analogous inequality, including
$q'(A)\le0$ at a right boundary optimum, gives
$a^*\le a-g_-/\mu$, while the proposed lower endpoint is at
most $a$.  Equality $a=a^*$ is immediate.  The unresolved-sign
interval has width $(g_+-g_-)/\mu$, proving
\eqref{eq:roundedprogress126}.  The transported cones of the
observation impose these same restrictions on their constrained
zeros; retaining further valid observations can only tighten them.

The floor of a nonnegative feasible midpoint is feasible, including
when $A$ is not a lattice point, and has the asserted error.  The
checked inequality therefore implies
$w_{k+1}\le w_k/2+h$.  Iterating this scalar inequality gives the
stated expression.  No contraction is inferred for a failed check,
for the first arbitrary actor center, or for a query that stops
before obtaining a valid derivative interval.  A lattice floor,
a terminal accuracy floor or the resource limit can leave the
stricter regret target unresolved.
\end{proof}

\subsection{What a finite direct-service return guarantees}

The frozen service takes $h=2^{-30}$, $M=7$, $\mu=2$ and
$\delta=15/2^{24}$, with the proved original $L$ appropriate to the
risk level.  Its ordinary tolerances are $\tau=2^{-8}$ and
$\omega=2^{-13}$; its separately reported stronger target is
$\tau_*=2^{-16}$.  It starts with the same paid actor or conventional
proposal as its comparator and buys its own first original support.
It does not inherit a completed comparison service without paying
for it.  The first actor-center departure is allowed without a
halving assertion; subsequent departures require the actual
bracket check in Lemma~\ref{lem:roundstop126}.  Every candidate is
the realized feasible action, including its rounding, rather than
an unimplemented ideal midpoint.

\begin{proposition}[Finite service qualification and the cost boundary]
\label{prop:directservice126}
Assume all completed support records satisfy the premises above and
all returned floating endpoints enclose their exact quantities.
Every successful direct-service return satisfies both ordinary
gates, either through Proposition~\ref{prop:directcertificate126}
or through its separately paid, independently executed qualified
fallback.  Each recorded distinct-action replacement has the strict
same-objective guarantee in Proposition~\ref{prop:strictchange126}.
Requalification of an unchanged action is not recorded as a policy
gain.  Reaching a resource or precision limit with an ordinarily
qualified incumbent does not imply attainment of $\tau_*$.

The direct route starts no more than $8\times7$ queries, never
purchases the same action at the same precision twice, and checks its cap
of 32,768 requested terminal contexts before starting each query.
A query at $n$ bins charges $2n+1$ requested contexts, with
$n\in\{32,64,128,256,512,1024,2048\}$.  These are finite budget
statements, not a bound on all arithmetic operations or elapsed time.
All failed or discarded direct work, actor acquisition and inference,
certificate work, and any completed fallback belong to the complete
service cost.
\end{proposition}

\begin{proof}
Before selection the service has no completed incumbent to improve.
Its first selection is restricted to candidates passing both exact
and outward ordinary gates.  After selection it reconsiders the
same action using all retained information.  Proposition
\ref{prop:directcertificate126} preserves exact qualification under
valid information accumulation.  A different action is selected
only if it also passes both gates and the strict replacement test;
Proposition~\ref{prop:strictchange126} proves the claimed gain.
Induction establishes qualification of every completed selected
policy.  The stronger stopping condition explicitly checks both
$G\le\tau_*$ and width.  Ending at a different declared limit
returns the qualified incumbent with its actual stronger-target
status.  A recoverable numerical failure returns a retained
completed qualified policy when one is available; otherwise the
service pays the entire qualified fallback.  If the fallback or an
infrastructure check aborts, there is no successful-return claim.

There are eight center rounds and seven declared bin levels per
round.  The visited-pair record prevents duplicate purchases.
Before beginning a new query the service adds its $2n+1$ requested
contexts only if doing so remains within the declared cap, and
retains the charge even if the query subsequently fails.  The
requested-context count is conservative about a failed partial
query and is distinct from the saved actual operation ledger.
The actor, sorting, terminal solving, support algebra, serialization,
and complete fallback can require additional work.  No part of the
budget argument licenses omitting those costs or assigning a zero
cost to a refused candidate.  The certificate and accounting rules
apply equally to neural, matched hinge, polynomial, spline, grid
and online proposals supplied with the same information.
\end{proof}
